Large Language Models (LLMs) are increasingly used as agents in social simulations, yet their ability to model complex interpersonal dynamics remains under-explored. This paper investigates whether LLMs capture the relational structure of persona-persona dislike, hypothesizing that personas separated by greater distance in a latent persona space will express stronger antipathy toward each other's preferences. Using the \dataset dataset and \gpt, we map personas into a multidimensional embedding space and conduct reciprocal preference evaluations. Our results demonstrate a statistically significant negative correlation ($r = -0.49$, $p < 0.001$) between persona distance and reciprocal preference agreement, indicating that LLMs inherently encode an "us-vs-them" logic based on value and lifestyle divergence. Furthermore, we introduce \schismogenesis, a task where the model adopts the "complete opposite" of a target persona. In this setting, the model demonstrates a remarkable 88\% rejection rate of the target's preferences, justifying its opposition through structural counter-identities. These findings suggest that LLMs do not merely model isolated traits but possess a sophisticated social map where identity is defined through relational opposition, with significant implications for polarization and cooperative AI design.
